% GENERATED FILE. Edit canonical lessons, then run npm run generate:books.
% canonical-source-hash: fnv1a64:ceabb6e6e963e580
% canonical-lessons: JA-C213-hikaru, JA-C213-tsuki, JA-C213-shizumu, JA-C213-noboru, JA-C213-atatamaru

\chapter{Sun, Moon and Stars}
\label{ch:ja-sun-moon-and-stars}
\hlchaptermodality{213}

\begin{chapteropening}
\textbf{By the end of this chapter:} \emph{I can talk about the sun, the moon and the stars.}

Complete the last lesson of chapter 213: I can talk about the sun, the moon and the stars.
\end{chapteropening}

\section[\texorpdfstring{\ja{ひかる} (hikaru)}{ひかる (hikaru)}]{\ja{ひかる} (hikaru) — to shine; to glitter}

\label{lesson:JA-C213-hikaru}



\begin{quote}
Before the new one: say the Japanese for to freeze, then the Japanese for to melt.
\end{quote}

\subsection*{You'll want to know: \ja{ひかる}}
\textbf{\ja{ひかる}} — \emph{hikaru} — \textquotedblleft{}to shine; to glitter\textquotedblright{}.

\begin{quote}
\textbf{\ja{こんや} \ja{ほしが} \ja{ひかります。}} — \emph{kon'ya hoshi ga hikarimasu} — \textquotedblleft{}The stars will shine tonight.\textquotedblright{}
\end{quote}

Polite \textbf{\ja{ひかります}}, \emph{hikarimasu}; a godan verb. The noun \textbf{\ja{ひかり}}, \emph{hikari}, means \textquotedblleft{}light\textquotedblright{}.

\begin{cousinweb}
\textbf{\ja{ひかる}} (\emph{hikaru}) is a native Japanese word (wago), inherited rather than borrowed from Chinese; its kanji was attached later as a kun'yomi, and its older origin is not settled (JMdict, KANJIDIC2).
\end{cousinweb}

\begin{grammarlens}[title={What you've built}]
One more word for everyday life at A2.
\end{grammarlens}

\subsection*{Your turn}
\begin{itemize}
\raggedright
  \item \emph{Say it:} \emph{hikaru}
  \item \emph{Say it:} \emph{hikaru}, once more
  \item \emph{Say it:} \emph{kon'ya hoshi ga hikarimasu} — The stars will shine tonight.
  \item \emph{Recall:} say the Japanese for to freeze, then the Japanese for to melt, then say \emph{hikaru} again
\end{itemize}

\begin{tcolorbox}[breakable,colback=teal!4,colframe=teal!35!black,title={Before you move on}]
What does \ja{ひかる} mean? (\textquotedblleft{}To shine; to glitter\textquotedblright{}.) Say it once more.
\end{tcolorbox}
\section[\texorpdfstring{\ja{つき} (tsuki)}{つき (tsuki)}]{\ja{つき} (tsuki) — the moon; a month}

\label{lesson:JA-C213-tsuki}



\begin{quote}
Before the new one: say the Japanese for to melt, then the Japanese for to shine.
\end{quote}

\subsection*{You'll want to know: \ja{つき}}
\textbf{\ja{つき}} — \emph{tsuki} — \textquotedblleft{}the moon; a month\textquotedblright{}.

\begin{quote}
\textbf{\ja{あしたは} \ja{よなかに} \ja{つきが} \ja{でます。}} — \emph{ashita wa yonaka ni tsuki ga demasu} — \textquotedblleft{}Tomorrow the moon will come out at midnight.\textquotedblright{}
\end{quote}

\begin{cousinweb}
\textbf{\ja{つき}} (\emph{tsuki}) is a native Japanese word (wago), inherited rather than borrowed from Chinese; its kanji was attached later as a kun'yomi, and its older origin is not settled (JMdict, KANJIDIC2).
\end{cousinweb}

\begin{grammarlens}[title={What you've built}]
One more word for everyday life at A2.
\end{grammarlens}

\subsection*{Your turn}
\begin{itemize}
\raggedright
  \item \emph{Say it:} \emph{tsuki}
  \item \emph{Say it:} \emph{tsuki}, once more
  \item \emph{Say it:} \emph{ashita wa yonaka ni tsuki ga demasu} — Tomorrow the moon will come out at midnight.
  \item \emph{Recall:} say the Japanese for to melt, then the Japanese for to shine, then say \emph{tsuki} again
\end{itemize}

\begin{tcolorbox}[breakable,colback=teal!4,colframe=teal!35!black,title={Before you move on}]
What does \ja{つき} mean? (\textquotedblleft{}The moon; a month\textquotedblright{}.) Say it once more.
\end{tcolorbox}
\section[\texorpdfstring{\ja{しずむ} (shizumu)}{しずむ (shizumu)}]{\ja{しずむ} (shizumu) — to sink; to set (of the sun)}

\label{lesson:JA-C213-shizumu}



\begin{quote}
Before the new one: say the Japanese for to shine, then the Japanese for the moon.
\end{quote}

\subsection*{You'll want to know: \ja{しずむ}}
\textbf{\ja{しずむ}} — \emph{shizumu} — \textquotedblleft{}to sink; to set (of the sun)\textquotedblright{}.

\begin{quote}
\textbf{\ja{たいようは} \ja{にしに} \ja{しずみます。}} — \emph{taiyō wa nishi ni shizumimasu} — \textquotedblleft{}The sun sets in the west.\textquotedblright{}
\end{quote}

Polite \textbf{\ja{しずみます}}, \emph{shizumimasu}; a godan verb.

\begin{cousinweb}
\textbf{\ja{しずむ}} (\emph{shizumu}) is a native Japanese word (wago), inherited rather than borrowed from Chinese; its kanji was attached later as a kun'yomi, and its older origin is not settled (JMdict, KANJIDIC2).
\end{cousinweb}

\begin{grammarlens}[title={What you've built}]
One more word for everyday life at A2.
\end{grammarlens}

\subsection*{Your turn}
\begin{itemize}
\raggedright
  \item \emph{Say it:} \emph{shizumu}
  \item \emph{Say it:} \emph{shizumu}, once more
  \item \emph{Say it:} \emph{taiyō wa nishi ni shizumimasu} — The sun sets in the west.
  \item \emph{Recall:} say the Japanese for to shine, then the Japanese for the moon, then say \emph{shizumu} again
\end{itemize}

\begin{tcolorbox}[breakable,colback=teal!4,colframe=teal!35!black,title={Before you move on}]
What does \ja{しずむ} mean? (\textquotedblleft{}To sink; to set (of the sun)\textquotedblright{}.) Say it once more.
\end{tcolorbox}
\section[\texorpdfstring{\ja{のぼる} (noboru)}{のぼる (noboru)}]{\ja{のぼる} (noboru) — to climb; to go up (a slope)}

\label{lesson:JA-C213-noboru}



\begin{quote}
Before the new one: say the Japanese for the moon, then the Japanese for to sink.
\end{quote}

\subsection*{You'll want to know: \ja{のぼる}}
\textbf{\ja{のぼる}} — \emph{noboru} — \textquotedblleft{}to climb; to go up (a slope)\textquotedblright{}.

\begin{quote}
\textbf{\ja{らいしゅう} \ja{やまに} \ja{のぼります。}} — \emph{raishū yama ni noborimasu} — \textquotedblleft{}Next week I will climb a mountain.\textquotedblright{}
\end{quote}

Polite \textbf{\ja{のぼります}}, \emph{noborimasu}; a godan verb.

\begin{cousinweb}
\textbf{\ja{のぼる}} (\emph{noboru}) is a native Japanese word (wago), inherited rather than borrowed from Chinese; its kanji was attached later as a kun'yomi, and its older origin is not settled (JMdict, KANJIDIC2).
\end{cousinweb}

\begin{grammarlens}[title={What you've built}]
One more word for everyday life at A2.
\end{grammarlens}

\subsection*{Your turn}
\begin{itemize}
\raggedright
  \item \emph{Say it:} \emph{noboru}
  \item \emph{Say it:} \emph{noboru}, once more
  \item \emph{Say it:} \emph{raishū yama ni noborimasu} — Next week I will climb a mountain.
  \item \emph{Recall:} say the Japanese for the moon, then the Japanese for to sink, then say \emph{noboru} again
\end{itemize}

\begin{tcolorbox}[breakable,colback=teal!4,colframe=teal!35!black,title={Before you move on}]
What does \ja{のぼる} mean? (\textquotedblleft{}To climb; to go up (a slope)\textquotedblright{}.) Say it once more.
\end{tcolorbox}
\section[\texorpdfstring{\ja{あたたまる} (atatamaru)}{あたたまる (atatamaru)}]{\ja{あたたまる} (atatamaru) — to get warm; to warm up (by itself)}

\label{lesson:JA-C213-atatamaru}



\begin{quote}
Before the new one: say the Japanese for to sink, then the Japanese for to climb.
\end{quote}

\subsection*{You'll want to know: \ja{あたたまる}}
\textbf{\ja{あたたまる}} — \emph{atatamaru} — \textquotedblleft{}to get warm; to warm up (by itself)\textquotedblright{}.

\begin{quote}
\textbf{\ja{ごごは} \ja{へやが} \ja{あたたまります。}} — \emph{gogo wa heya ga atatamarimasu} — \textquotedblleft{}The room gets warm in the afternoon.\textquotedblright{}
\end{quote}

Polite \textbf{\ja{あたたまります}}, \emph{atatamarimasu}; a godan verb.

\begin{cousinweb}
\textbf{\ja{あたたまる}} (\emph{atatamaru}) is a native Japanese word (wago), inherited rather than borrowed and built on the same stem as \emph{atatakai}, \textquotedblleft{}warm\textquotedblright{}; its kanji was attached later as a kun'yomi (JMdict, KANJIDIC2).
\end{cousinweb}

\begin{grammarlens}[title={What you've built}]
One more word for everyday life at A2.
\end{grammarlens}

\subsection*{Your turn}
\begin{itemize}
\raggedright
  \item \emph{Say it:} \emph{atatamaru}
  \item \emph{Say it:} \emph{atatamaru}, once more
  \item \emph{Say it:} \emph{gogo wa heya ga atatamarimasu} — The room gets warm in the afternoon.
  \item \emph{Recall:} say the Japanese for to sink, then the Japanese for to climb, then say \emph{atatamaru} again
\end{itemize}

\begin{tcolorbox}[breakable,colback=teal!4,colframe=teal!35!black,title={Before you move on}]
What does \ja{あたたまる} mean? (\textquotedblleft{}To get warm; to warm up (by itself)\textquotedblright{}.) Say it once more.
\end{tcolorbox}
